\documentclass[11pt,a4paper]{article}
\usepackage[UTF8,fontset=none]{ctex}
\setCJKmainfont{Noto Serif CJK SC}[BoldFont={Noto Sans CJK SC Bold}]
\setCJKsansfont{Noto Sans CJK SC}
\setCJKmonofont{Noto Sans Mono CJK SC}
\setmainfont{Liberation Serif}
\setsansfont{Liberation Sans}
\setmonofont{Liberation Mono}
\usepackage{amsmath,amssymb,graphicx,booktabs,array}
\usepackage{geometry,caption,hyperref,flafter}
\geometry{left=22mm,right=22mm,top=23mm,bottom=24mm}
\hypersetup{colorlinks=true,linkcolor=black,citecolor=black,urlcolor=blue,
 pdftitle={Category Priors and Geometric Feedback for Budget-Constrained Active Observation of Facilities},pdfauthor={Zhaoyu Li}}
\captionsetup{font=small,labelfont=bf,labelsep=quad}
\setlength{\parskip}{3pt}
\setlength{\emergencystretch}{3em}
\renewcommand{\topfraction}{0.9}
\renewcommand{\bottomfraction}{0.85}
\renewcommand{\textfraction}{0.07}
\renewcommand{\floatpagefraction}{0.8}
\linespread{1.13}
\urlstyle{same}
\Urlmuskip=0mu plus 2mu
\providecommand{\tightlist}{\setlength{\itemsep}{0pt}\setlength{\parskip}{0pt}}
\providecommand{\passthrough}[1]{#1}
\providecommand{\pandocbounded}[1]{#1}
\title{\textbf{Category Priors and Geometric Feedback for Budget-Constrained Active Observation of Facilities}}
\author{Zhaoyu Li}
\date{28 September 2026\quad Review draft}
\DeclareMathSizes{10.95}{10.95}{8}{7}
\clubpenalty=10000
\widowpenalty=10000
\displaywidowpenalty=10000
\newlength{\ReviewTableWidth}
\renewcommand{\abstractname}{Abstract}
\renewcommand{\figurename}{Figure}
\renewcommand{\tablename}{Table}
\begin{document}
\maketitle
\begin{abstract}
After changes to an industrial layout, robots need current observations of the surrounding space and equipment. Observing a region, however, does not ensure that its facilities have been adequately reconstructed. For facility documentation under a limited motion budget, this study combines category-conditioned configuration beliefs, predicted observation value, and feedback from measured geometry. A hierarchical decision and local execution architecture connects four module interfaces, with atomic decisions in the local implementation and observation macros in the separate multi-facility extension. Public templates support belief updates and planning, while the reconstructed surface contains only fused measured depth. Virtual validation uses separate static mapping tasks. A paired confirmation experiment uses two previously seen parent layouts, each with two configurations. Relative to an active geometric control with the same planner, category information improves the mean joint documentation score by 6.1638\%, with two wins and two ties across four pairs and identical mean coverage. In a separate development experiment with incorrect category information, the complete correction policy, including measured correction and anticipated future correction, improves the mean joint score by 2.6241\%. A further fixed 128-trial matrix retains 96 qualified completions and 32 low-budget planning failures. Category gains at budget 42 are 6.25\% on original layouts and 4.23\% on same-family variants, while budget 54 produces a −0.70\% difference on original layouts. A separate four-facility study comprises three controller versions and 36 attempts; shared-reliability and fixed category-prior outcomes tie in each of the three layouts for every version. The local results support a conditional benefit for early observation selection; they do not establish an additional benefit from cross-instance sharing.
\end{abstract}
\noindent\textbf{Keywords:} facility documentation; active observation; category priors; geometric feedback; budget-constrained planning

\hypertarget{introduction}{%
\section{Introduction}\label{introduction}}

Embodied robots connect perception, decision making, and physical action. Their applications in manufacturing and logistics are expanding, although scaling industrial deployments remains challenging {[}9{]}. Equipment rearrangement, workstation adjustment, and changes in material placement can make an existing environmental representation unsuitable for a new task. This study considers how a robot can actively acquire current spatial and equipment geometry after such a reconfiguration. The environment remains static during each mapping task; tracking moving obstacles and dynamic obstacle avoidance are outside the task.

Active SLAM uses robot motion selection to improve environmental modeling {[}10{]}. The present study addresses its observation-planning and mapping-quality aspects. Industrial facility documentation requires both spatial coverage and observations of equipment surfaces. A robot may cover an accessible corridor while leaving rear or recessed surfaces poorly reconstructed. Additional observations of open space do not necessarily recover these surfaces. With a narrow-field-of-view depth camera, the direction of approach and the timing of turns determine how much useful surface evidence can be acquired within a limited motion budget.

Category knowledge can provide a directional cue. When equipment categories are associated with service sides or openings, a robot may select a useful observation side before the relevant structure becomes geometrically distinguishable. However, categories can be incorrect, and geometric planners can actively acquire diagnostic observations and change direction. A higher score along a semantically guided trajectory therefore does not, by itself, isolate a semantic effect. This study asks a specific question: with shared templates, a safe pose graph, a planner, and a reconstruction backend, can category information reduce the cost of diagnosing a configuration and subsequently redirecting observation?

SWAP connects semantic inspection to observation requirements {[}2{]}, and VISTA combines task relevance with directional coverage {[}3{]}. For budget-constrained documentation of industrial facilities, this study implements a category-prior--geometric-feedback loop, connects observation decisions to execution checks, and isolates the contribution of category information through a shared active geometric control. The contribution lies in the executable connection between category cues, actions, and reconstructed surfaces; Bayesian updating and information value follow established principles. Controlled virtual validation is the primary evidence, with physical implementation left for supplementary demonstration.

\hypertarget{related-work}{%
\section{Related Work}\label{related-work}}

Autonomous exploration coordinates environmental representations, target selection, and local motion. A survey of active SLAM connects motion decisions with map estimation and discusses modular and belief-space planning formulations {[}10{]}. Active Neural SLAM is an early example using learned mapping and hierarchical policies {[}1{]}, while TARE coordinates exploration through detailed local and sparse global representations {[}5{]}. The present study uses hierarchical decisions and local execution to allocate facility observations within a common safe graph and action budget. Its internal geometric control can revisit locations, seek configuration-disambiguating evidence, and replan. The controlled variable is the contribution of category information beyond these shared active planning capabilities.

Active reconstruction considers where the next observation should be acquired. GenNBV uses reinforcement learning with geometric, semantic, and action representations to select viewpoints in a five-dimensional free space and studies cross-scene next-best-view behavior {[}6{]}. Facility documentation also involves self-occlusion and insufficient viewing directions. Here, the task is restricted to a finite ambiguity between public configurations, discrete ground motion, and explicit budget constraints. All internal controls use the same Open3D depth-fusion backend {[}4{]}. Templates inform prediction, whereas only executed observations contribute to reconstruction, separating observation allocation from surface completion by prior knowledge.

Semantic inspection connects the choice of targets to observation sufficiency. SWAP combines volumetric exploration, semantic target completion, and surface inspection with resolution and incidence-angle requirements {[}2{]}. Semantics-Aware Predictive Inspection Path Planning (SPP) exploits repeated structure in semantic scene graphs to predict unseen targets in industrial inspection {[}7{]}. Accordingly, the present motivation is more specific than proposing semantic revisiting or structural prediction: when current geometry cannot distinguish observation directions, can category--structure associations reduce the cost of choosing a direction, and can measured geometry correct a misleading association?

Recent semantic active mapping also combines task relevance and observation value. VISTA uses online semantic Gaussian maps for open-vocabulary tasks and scores receding-horizon trajectories using semantic relevance and directional coverage {[}3{]}. ActiveSGM combines semantic and geometric uncertainty to predict the informativeness of candidate observations {[}8{]}. This study instead implements category-conditioned configuration weights and measured-residual feedback on a CPU. Its SWAP-I and VISTA-I comparisons adapt selected planning mechanisms to common candidates and fusion conditions; they are not reproductions of the complete original systems. Paired controls, reconstructed endpoints, and recorded belief--action changes together establish the scope of the resulting application evidence.

\hypertarget{task-and-evaluation}{%
\section{Task and Evaluation}\label{task-and-evaluation}}

The facility configuration remains fixed within each task, and different configurations define separate mapping tasks. Layout reconfiguration motivates the application; the experiments do not estimate temporal changes or incrementally update a historical map. The discrete robot pose is \(s_t=(x_t,y_t,\theta_t)\), and the facility configuration is \(h\in\{0,1\}\). The robot receives two public templates, a facility coordinate frame, and a common safe pose graph. It does not know the true configuration. Runtime inputs comprise RGB-D, scan, and pose observations already acquired through execution. The category--structure relation is manually specified public knowledge. Evaluation ground truth and actual observations from unexecuted actions are unavailable to the controller.

A forward motion of 1 m or a turn of 90° costs one action. The total budget is B actions: an 18-action shared observation prefix is followed by at most B−18 autonomously selected actions. The original confirmation uses 42 actions; the budget extension uses 30, 42, and 54. A valid task requires at least 80\% coverage, no collision, completion within budget, and return to the initial position and orientation. Thus, the experiment evaluates documentation decisions after shared initial observation rather than unknown-environment exploration from action zero.

The evaluator uses a 0.2 m grid and a 0.2 m robot radius to derive the start-reachable safe-cell set \(\mathcal R\) from scene ground truth. Its size is 588 for P00 and 668 for P01. This set is used only after outputs are sealed and is not an additional ground-truth input to the controller. Coverage and the primary score are

\[
C_{\rm map}=\frac{|\{c\in\mathcal R:m_T(c)\ne-1\}|}{|\mathcal R|},
\qquad F_{1,\tau}=\frac{2P_\tau R_\tau}{P_\tau+R_\tau},
\qquad J_5=C_{\rm map}F_{1,0.05}.
\]

Unknown cells have value −1; cells marked occupied also count as known. Coverage consequently measures observation status rather than occupancy classification accuracy. Surface precision, recall, and their harmonic mean follow the distance-threshold evaluation principle of Tanks and Temples {[}14{]}. The common facility region and the product with coverage are task-specific definitions. Surface precision samples the predicted mesh and queries the complete reference surface; recall samples fixed exterior vertical reference surfaces and queries the prediction. Before evaluation, predictions are cropped to the common bounding box of both templates, expanded by 0.20 m, and the known ground is removed. No hole filling or registration is applied. Errors outside this shared region do not enter precision, so the score describes the facility evaluation domain. Thresholds are fixed at 2, 5, and 10 cm, with 5 cm primary. Task validity is reported separately from quality; an empty prediction receives zero quality, and \(F_{1,\tau}=0\) when \(P_\tau+R_\tau=0\).

The intended decision objective is to maximize expected measured documentation quality over policies using only permitted, already acquired information:

\[
\max_{\pi\in\Pi_{\rm paid}}\mathbb E[J_5(M_T,G^\star)],
\qquad \sum_{t=0}^{T-1}c(a_t)\le B,\quad s_T=s_0.
\]

Here \(G^\star\) denotes the evaluator's reference geometry, not an online input. The coverage and collision requirements above further determine task validity. Both G and S may use observed geometric deficiencies, acquire diagnostic measurements, and revisit locations; S additionally conditions decisions on the acquired category cue. Maximizing a template-exposure proxy is an approximation to this objective, and cannot guarantee an increase in measured \(J_5\).

Figure 1 relates the two public configuration hypotheses to occlusion and the saved routes. Both zero-gain configurations are shown alongside the conditions in which category information changes the observation side.

\begin{center}
\begin{minipage}{\linewidth}
\centering
\includegraphics[width=\linewidth,height=.72\textheight,keepaspectratio]{docs/thesis/figures/article_trajectories_20260928/local_paired_trajectories.pdf}
\captionof{figure}{All four V36 confirmation pairs with saved routes and complete primitive-action strips. Panels share a metric scale; P01 is displayed in its declared equipment frame. Grey marks the common 18-action prefix. Open circles denote in-place turn positions; arrows show sampled recorded headings. Open diamonds indicate the first nonzero applied geometric residual, and filled diamonds the shared start/return. Both h0 trajectories coincide; both h1 pairs first diverge at executed action 19. Every action obtains a saved sensor frame; these runs have no extra observe primitive. J5 is copied from sealed endpoints.}\label{fig:writing-1}
\end{minipage}
\end{center}

\hypertarget{category-priors-and-geometric-feedback}{%
\section{Category Priors and Geometric Feedback}\label{category-priors-and-geometric-feedback}}

\hypertarget{module-interfaces-and-executed-controller}{%
\subsection{Module Interfaces and Executed Controller}\label{module-interfaces-and-executed-controller}}

The system organizes observation, planning, and feedback through hierarchical decisions and local execution. Table 1 describes the four CPU interfaces. The decision layer selects a pose subgoal from the configuration belief and remaining budget; the execution layer checks action legality and the return budget before new measurements update the state. In the local binary implementation, each subgoal corresponds to the next atomic action; Section 4.6 defines a separate scene-level extension. Controlled category inputs and deterministic execution guards complete the loop. Experiments isolate category information and the complete correction policy, rather than assign an independent performance contribution to each interface.

\par\noindent\begin{minipage}{\linewidth}
\textbf{Table 1. Responsibilities of the four CPU module interfaces.}

\begin{center}
\small
\setlength{\tabcolsep}{4pt}
\setlength{\ReviewTableWidth}{\dimexpr\linewidth-4\tabcolsep\relax}
\begin{tabular}{@{}>{\raggedright\arraybackslash}p{0.28\ReviewTableWidth} >{\raggedright\arraybackslash}p{0.36\ReviewTableWidth} >{\raggedright\arraybackslash}p{0.36\ReviewTableWidth}@{}}
\toprule
Interface & Inputs & Processing and output \\
\midrule
OV-SDF & Executed RGB, pose, and measured-map summary & Register category support; output category log-prior and observation records \\
STGHP & Common graph, template exposure, configuration weights, remaining budget & Compare feasible continuations; output the next pose and candidate values \\
RPN-UQ & Selected edge, remaining budget, return distance & Output an execution permission or rejection record \\
IGCR & Measured depth/scan and public predictions at a new pose & Compare residuals; output a geometric evidence increment and revised weights \\
\bottomrule
\end{tabular}
\end{center}
\end{minipage}


\begin{center}
\begin{minipage}{\linewidth}
\centering
\includegraphics[width=\linewidth,height=.72\textheight,keepaspectratio]{docs/thesis/figures/system_architecture_20260928/system_architecture.pdf}
\captionof{figure}{Platform interfaces, the local binary implementation of the four-module CPU loop, and separate mapping and evaluation. OV-SDF and IGCR jointly update configuration belief; STGHP selects the next atomic action for RPN-UQ to check before execution and new observation. Public templates support belief and planning, while TSDF fuses measured depth and ground truth enters only evaluation of sealed outputs. The robot is a generative composite based on the user-provided platform photograph for interface illustration; sensor, map, and mesh thumbnails are schematic.}\label{fig:writing-2}
\end{minipage}
\end{center}

\hypertarget{measured-belief-state}{%
\subsection{Measured Belief State}\label{measured-belief-state}}

Inputs comprise the common safe pose graph \(\mathcal G\), initial pose \(s_0\), public templates \(T_0,T_1\), budget \(B\), and shared prefix. Both methods receive template predictions of depth, scans, and potential exposure at each pose. Outputs include executed actions, observation records, and a measured map. Independent endpoint evaluation follows output sealing.

After action \(t\), the planning state is

\[
x_t=(s_t,b_t,M_{0,t},M_{1,t},p_t,\mathcal V_t),\qquad b_t=B-t,
\]

where \(M_{h,t}\) is the union of template exposures associated with previously observed poses, and \(\mathcal V_t\) is the visited-pose set. With \(p_t\) denoting the weight assigned to \(h=0\), the belief update is

\[
L_t^g=\operatorname{clip}_{[-24,24]}(L_{t-1}^g+\delta_t),
\qquad p_t=\sigma(L_t^g+L_t^s).
\]

At a newly visited pose, \(\delta_t\) is the mean truncated-loss difference \(e_1-e_0\) over valid depth and scan modalities, clipped to \([-6,6]\). It is zero for repeated poses or disabled feedback. Depth and scan require at least eight valid pixels and three valid beams, respectively. Each loss is half the squared normalized residual, capped at 12.5. Category registration requires at least 16 exact-color pixels at each of two distinct XY positions, with no competing marker color in those frames. Registration sets \(L_t^s=\pm\log9\); repeated frames do not multiply the prior, and conflicting registered categories clear it. The geometric control G always sets \(L_t^s=0\). These fixed-scale pseudo-likelihood weights are not calibrated probabilities.

Residuals are evaluated only where the template predictions differ. Measured zero depth and zero scan range are treated as missing. Depth residuals use the reference disparity domain, \((57.6/z-57.6/\widehat z_h)/0.25\), while scan residuals use \((r-\widehat r_h)/0.02\). A missing template depth return incurs the capped loss. Without a valid modality, no geometric evidence is added. The scan scale of 0.02 m controls evidence weighting and is not an imposed simulator noise level.

An incorrect prior has a finite correction threshold. If \(L^s>0\) favors configuration 0, assigning configuration 1 a support weight of at least \(\eta>1/2\) requires

\[
L_t^g\le-L^s-\operatorname{logit}(\eta).
\]

If every eligible counter-observation contributes at most \(-\epsilon<0\), a sufficient count from \(L_0^g\) is \(\max\{0,\lceil(L_0^g+L^s+\operatorname{logit}(\eta))/\epsilon\rceil\}\), provided the cumulative cap permits the threshold. With the present cap this requires \(24\ge L^s+\operatorname{logit}(\eta)\). This is an algebraic statement about the implemented support weights, not a calibrated probability or a guarantee that informative measurements will be obtained. First-pose counting suppresses exact repeats but does not establish independence between nearby views. A measurement can therefore remain useful for depth fusion without qualifying as additional belief evidence.

\hypertarget{limited-diagnostic-lookahead-and-budgeted-execution}{%
\subsection{Limited Diagnostic Lookahead and Budgeted Execution}\label{limited-diagnostic-lookahead-and-budgeted-execution}}

The predictive objective is \(\widetilde J_h(M_h)=\widetilde C_h(M_h)\widetilde F_h(M_h)\), using ideal template ground exposure and the fraction of potential exterior vertical surface area. It is distinct from measured \(J_5\). The predicted stopping value is

\[
Z(s,M,p)=
\begin{cases}
p\widetilde J_0(M_0)+(1-p)\widetilde J_1(M_1),
&s=s_0,\ \min_h\widetilde C_h(M_h)\ge0.8,\\
-\infty,&\text{otherwise}.
\end{cases}
\]

Belief-space decision making connects observations to subsequent actions {[}12{]}, while informative path planning jointly considers information quality and resource budgets {[}13{]}. Here this connection is approximated by one future diagnostic event per planning call. The dynamic-programming state includes \(q\in\{0,1\}\), indicating whether one future diagnostic event remains available in the current lookahead. Action \(a\) moves to \(s'\), updates exposure to \(M'_h=M_h\cup E_h(s')\), and consumes one unit of budget. If \(q=1\) and \(s'\) is an unobserved diagnostic node, the prediction branches using an assumed diagnostic accuracy \(r=0.99\):

\[
\lambda_0=pr+(1-p)(1-r),\quad \lambda_1=1-\lambda_0,
\qquad p^{(0)}=\frac{pr}{\lambda_0},\quad
p^{(1)}=\frac{p(1-r)}{\lambda_1}.
\]

\[
Q_q(a)=
\begin{cases}
\sum_{y=0}^{1}\lambda_y V_0(s',b-1,M',p^{(y)}),&\text{diagnostic event},\\
V_q(s',b-1,M',p),&\text{otherwise},
\end{cases}
\qquad V_q=\max\{Z,\max_a Q_q(a)\}.
\]

Previously measured poses are excluded from diagnostic candidates at the beginning of each planning call. Diagnostic nodes are identified by comparing the two public templates: signed evidence under their respective predictions must reach the corresponding \(\pm\log99\) thresholds. No further information event is expanded after the forecast diagnosis. If either diagnostic branch is infeasible, the action is infeasible. A state also has value \(-\infty\) when its return distance exceeds its remaining budget. Only the first selected action is executed before replanning from a new measurement; forecast weights never overwrite the measured belief record.

The measured map and template exposure sets are maintained separately. The map fuses acquired depth, whereas \(M_h\) describes exposure expected if template \(h\) is correct; template surfaces are never inserted into the mesh. The shared 18-action prefix returns to the initial pose while updating beliefs. Only subsequent actions depend on the resulting weights.

Execution must satisfy \(1+d_{\rm return}(s')\le b_t\), with return distance measured to the complete initial pose. On a fixed legal graph with correctly modeled action costs, this preserves \(d_{\rm return}(s_t)\le b_t\) by induction: after an allowed action, the reserved shortest return path costs no more than the remaining budget. The same certificate extends to a macro-action of cost \(c\) ending at \(u\) if \(c+d_{\rm return}(u)\le b_t\); every intermediate prefix retains a return path through the unexecuted suffix. All observation actions must be charged. This establishes existence of an affordable return path under the graph model, not completion of that path after a collision or computational failure. The reported controller executes one atomic action at a time. Voluntary stopping and forced termination are recorded separately, and termination alone does not establish task validity. Algorithm 1 summarizes the complete execution loop.

\par\noindent\begin{minipage}{\linewidth}
\small
\textbf{Algorithm 1. Receding-horizon planning with budgeted observations.}

\begin{verbatim}
Input: common safe graph and two templates, budget B, shared prefix pre,
       mode G/S/X/Xnf
Output: executed actions and observations, measured map, stop reason
Initialize: Lg=Ls=0; M0=M1={}; visited={}; t=0
Repeat:
  Read the initial frame or the frame from the latest action
  Fuse that frame exactly once through the execution driver
  Form observations without the true configuration identifier
  Verify action, step index, and pose
  OV-SDF: register category and map summaries
  IGCR: update Lg from measurements at a newly visited pose
  Add the current pose's template exposures to M0 and M1
  visited <- union(visited, {s})
  If collision has occurred or t=B:
    record the forced termination reason and exit
  If t<18: a <- pre[t]
  Otherwise:
    D <- public-template diagnostic nodes minus visited
    STGHP: solve Vq; use q=0 for Xnf and q=1 otherwise
    Report infeasibility or resource-limit failure as failure
    If the stopping value is optimal within 10^(-12):
      record voluntary stopping and exit
    a <- first optimal action within tolerance in fixed graph-edge order
  RPN-UQ: check edge legality and 1+dreturn(s') <= B-t
  If rejected: record the termination reason and exit
  Execute a; t <- t+1
Seal outputs; independently evaluate coverage, surface F1, J5, and validity
\end{verbatim}
\end{minipage}


Mode X exchanges the category interpretation. Xnf additionally disables both measured geometric updates and future diagnostic anticipation; X−Xnf therefore evaluates a complete correction policy. Each planning call is capped at 1.5 million cached states and 20 s. Time is checked every 1024 newly cached states and at selection completion, so the threshold is not a hard real-time bound.

For \(N\) poses, budget \(B\), and exposure bitsets of lengths \(m_0,m_1\), the memoized state count has the conservative bound \(S\le3N(B+1)2^{m_0+m_1}\): one initial weight and at most two post-diagnosis weights are considered within a call. With maximum graph out-degree \(d\), bitset-update cost \(T_m\), and terminal-evaluation cost \(T_z\), planning costs \(O(S(dT_m+T_z))\), in addition to return-distance preprocessing. The exposure dependence is exponential; caching and resource caps do not establish polynomial scaling or global optimality for measured reconstruction quality.

G and the category-conditioned method S share the same predictive belief planner. Both can actively diagnose a configuration and change observation direction; S additionally uses the category prior. They share an Open3D CPU implementation {[}4{]} of volumetric range integration {[}11{]}, with a 4 cm voxel size and 12 cm truncation distance. Only measured depth is fused. Actions, observations, and meshes are saved before endpoint evaluation.

\hypertarget{fixed-settings-and-a-recorded-two-configuration-example}{%
\subsection{Fixed Settings and a Recorded Two-Configuration Example}\label{fixed-settings-and-a-recorded-two-configuration-example}}

The category strength of 0.9 and forecast accuracy of 0.99 are fixed design assumptions, not measured recognition accuracy or calibrated sensor probabilities. The former supplies a finite preference that contrary geometry can overturn; the latter retains a nonperfect future diagnosis. The loss cap of 12.5, per-pose evidence limit of 6, and cumulative limit of 24 bound extreme residual influence. Repeated poses do not accumulate geometric evidence again, although new depth still enters TSDF fusion. A cache reuses dynamic-programming subproblems within each call, with the stated state/time limits; equal-valued actions use the fixed graph-edge order and \(10^{-12}\) tolerance. These choices are shared across controls and are not claimed to be optimal parameter values.

A saved P00 development example makes the connection to action explicit. Let \(w_t=1-p_t\) be the weight of configuration 1. After the shared prefix, G has \(w_{18}=0.5\): the left and right continuation values are both 0.650944, so the fixed tie rule selects left. For configuration 1, S has \(w_{18}=0.9\); the left/right values become 0.668008/0.711865 and action 19 turns right. For configuration 0 the category preference and value ordering reverse, and S turns left, coinciding with G. This explains why both a positive effect and a tie can occur in the same layout. All values are recorded terminal planning proxies, not measured F1 or immediate observation rewards.

With the misleading prior X in configuration 1, \(w_{18}=0.1\). At step 28, 1092 informative depth pixels produce mean clipped losses of 12.205313 and 0.505694 for the two templates. The applied log-evidence is −6, giving

\[
w_{28}=1-\sigma(-6+\log9)=0.978178.
\]

The forward/left/right continuation values are 0.574541/0.554660/0.555707, so action 29 moves forward. Its return cost is 13 actions with 14 remaining, satisfying \(1+13\le14\). At step 41 only one action remains; the robot is at the start position and turns to restore the starting orientation. Xnf retains weight 0.1 and turns right at action 29. Because Xnf also disables future diagnosis, its endpoint contrast estimates the whole correction policy; Figure 8 shows the associated measured belief timeline. This example uses saved development logs and adds no endpoint samples.

\hypertarget{when-can-a-category-cue-save-diagnostic-cost}{%
\subsection{When Can a Category Cue Save Diagnostic Cost?}\label{when-can-a-category-cue-save-diagnostic-cost}}

A simplified decision model clarifies the proposed mechanism. Assume equally likely configurations that the common initial geometry cannot distinguish. Two return-feasible routes yield utility \(U_+\) when matched to the configuration and \(U_-\) otherwise, with \(\Delta=U_+-U_->0\). A category cue available in the shared prefix has symmetric reliability \(r_c\ge1/2\). Both policies can instead acquire a symmetric geometric diagnosis of reliability \(r_d\ge1/2\), conditionally independent of the cue given the configuration. Both routes remain feasible after diagnosis, whose opportunity loss is \(\kappa\ge0\), expressed in utility units. These are analytical assumptions; \(r_c\) and \(r_d\) are not the implementation's fixed prior and forecast settings.

For a current configuration probability \(p\), direct selection is correct with probability \(m=\max(p,1-p)\). Selecting the better route after each possible diagnosis gives

\[
A(p,r_d)=\max\{pr_d,(1-p)(1-r_d)\}
+\max\{p(1-r_d),(1-p)r_d\}
=\max\{m,r_d\}.
\]

The net diagnostic value is therefore \(\Delta\max(0,r_d-m)-\kappa\). Allowing each policy to choose whether to diagnose yields

\[
V_G^*=U_-+\max\{\Delta/2,r_d\Delta-\kappa\},
\] \[
V_S^*=U_-+\max\{r_c\Delta,\max(r_c,r_d)\Delta-\kappa\}.
\]

Thus G benefits from diagnosis when \(\kappa<(r_d-1/2)\Delta\), whereas S requires \(\kappa<(r_d-r_c)_+\Delta\). For illustration only, \(U_-=0.50\), \(U_+=0.90\), \(r_c=0.80\), \(r_d=0.95\), and \(\kappa=0.08\) give \(V_G^*=0.80\) through diagnosis and \(V_S^*=0.82\) through direct category-based selection. With \(\kappa=0.02\), both diagnose and attain 0.86. These values are synthetic calculations, not reconstruction measurements.

The advantage requires an informative cue, consequential route choice, and a cost of resolving ambiguity. It can disappear when geometry already resolves the configuration, routes have equal value, or diagnosis is cheap. If diagnosis removes a route from the feasible set, that set must be re-evaluated rather than applying the closed-form result. The argument uses established belief-based decision principles {[}12,13{]}; it neither introduces a new Bayesian rule nor guarantees a TSDF improvement from additional views.

\hypertarget{independent-scene-level-extension}{%
\subsection{Independent Scene-Level Extension}\label{independent-scene-level-extension}}

A separate implementation extends local binary observation to four facilities in two repeated categories, each with planar, recessed, louvered, or open-frame structure. Policy labels are specific to each implementation: local S enables category conditioning, whereas scene-level B already includes that information and scene-level S additionally shares class-model evidence. Their effect sizes are not pooled. Instances and label planes arise from measured support; true facility poses and structures remain private. OV-SDF maintains associations and structure beliefs, IGCR supplies bounded depth evidence, STGHP selects observation macros, and RPN-UQ checks primitive execution and full-pose return. The supplied navigation graph, exact poses, synthetic labels, and nominal structure family remain controlled assumptions. TSDF fuses measured depth only. Within each scene family, development and reserved test layouts share the fixed category--configuration composition; the reserved layouts can test spatial transfer, not a new semantic-relation distribution.

For instance \(i\), cumulative geometric evidence \(L_i\) combines with a mixture of uniform and category priors:

\[
b_i(h)\propto[(1-\rho_i)q_0(h)+\rho_iq_1(h\mid c_i)]e^{L_i(h)}.
\]

G uses the uniform prior; B fixes \(\rho_i=1/2\); S updates the model weight from qualified same-category peers. Leave-one-out sharing applies own-instance evidence once and replaces cumulative peer messages. Category qualification requires two eligible frames, each containing at least four matching marker pixels, without conflict or association ambiguity. All methods retain measured geometry and the same acquisition rules. G, B, and S can purchase a diagnostic view; NBV uses G's belief with myopic direct-view scoring. Thus S/B tests sharing, B/G tests category conditioning, and G/NBV tests diagnostic lookahead.

At most 32 poses are compared, using expected nominal unseen area and an unknown-space proxy divided by arrival, observation, and return costs. Diagnostic alternatives consider one paid view and one subsequent common pose; instance utilities are summed before selecting that pose. Only the first macro is committed. Translation by 0.25 m, rotation by 30°, and explicit observation each cost one action and acquire a frame. One initial frame is free; geometry-only initialization consumes at most \(\lfloor B/5\rfloor\) actions. This replaces the old fixed prefix and binary dynamic program. Equations A1--A6 in the \href{/root/NSO/docs/thesis/ARTICLE_SUPPLEMENTARY_METHODS_20260928.pdf}{Supplementary Methods} specify belief transfer, macro utility, and limitations.

The extension evaluates measured-free coverage \(C_{\rm nav}\) on a fixed start-connected robot-center grid and \(Q=\tfrac14\sum_iF_{1,i}\) over all four facilities, including undiscovered ones. The joint score is \(J_{\rm nav}=C_{\rm nav}Q\). Complete predicted meshes are evaluated against fixed observable references at 5 cm, without prediction-ROI cropping or semantic weighting. These definitions differ from Section 3; the scores are not pooled.

The sharing model has a bounded effect. With one eligible peer and equal initial model weights, its evidence ratio satisfies \(r\in[\min_h q_1(h)/q_0(h),\max_h q_1(h)/q_0(h)]\), and \(\rho=r/(1+r)\). This bounds the transferred prior; an instance's own evidence can still concentrate its posterior. For fixed own evidence, each instance's contribution to a candidate-pair contrast is a linear-fractional function of \(\rho_i\) with a positive denominator, whose extrema occur at the interval endpoints. Summing per-instance extrema gives a conservative bound on a possible ranking change. This bound concerns a fixed direct-view candidate pool. It does not bound diagnostic replanning, future candidate generation, or measured reconstruction quality.

Two common-component revisions are evaluated separately. The first removes a verified measured floor band from instance-association candidates, while retaining full depth for occupancy and TSDF. The second discounts nominal surface opportunities already exposed at the same optical center. For nominal point \(x\), candidate visibility \(V_a(x)\), distance \(d(x,\mathcal P_t)\) from measured support, and accepted historical visibility masks \(V_j(x)\), its eligible area is

\[
A_t(a)=\sum_x a_x V_a(x)\,\mathbf 1[d(x,\mathcal P_t)>0.08\,\mathrm m]\,
\mathbf 1[x\notin\bigcup_{j\in\mathcal H_t(a)}V_j].
\]

Here \(\mathcal H_t(a)\) contains only already acquired frames with matching optical center, intrinsics, and image shape. The common initial frame is included; other history is paid acquisition. New viewing directions can retain unexposed nominal area, and observations from other optical centers are not excluded by this rule. Exclusion prevents repeatedly rewarding the same attempted opportunity under differently scaled templates; it neither certifies that depth was obtained nor adds any surface to the reconstruction. Consequently, recovery from missing depth and benefits from repeated measurement are not modeled by this discount. The original, ground-aware, and exposure-discounted controllers remain separate versions with complete outcomes.

\hypertarget{virtual-experiments-and-results}{%
\section{Virtual Experiments and Results}\label{virtual-experiments-and-results}}

\hypertarget{comparisons-and-common-settings}{%
\subsection{Comparisons and Common Settings}\label{comparisons-and-common-settings}}

The camera produces 96×72 images with a 4 m range. The planar scanner covers 360° with 180 beams and an 8 m range. Simulated poses are exact. Depth includes numerical noise with a standard deviation of 0.25 pixels in the reference disparity domain; scans have no additional random noise. These settings provide a controlled comparison rather than a calibrated model of physical sensors, and do not evaluate SLAM localization accuracy.

The confirmation experiment freezes the controller and uses a preselected held-out noise seed. Two parent layouts previously involved in development each contain two configurations, yielding eight G/S runs and four matched pairs. Facilities remain static within every run. Independent online replays reproduce all results but do not increase the sample count. Three comparisons address different questions: G/S isolates category availability, the incorrect-prior development comparison evaluates the complete correction policy, and the external CPU adapters compare observation-allocation mechanisms with different objectives and planning details.

The archived preparation manifests for development, confirmation, and CPU mechanism comparisons record Python 3.12.3, NumPy 1.26.4, SciPy 1.11.4, and Open3D 0.19.0. OpenMP, OpenBLAS, and MKL thread counts are fixed to one. The historical manifests do not record the CPU model or core count, so the reported timing is restricted to the recorded implementation conditions and is not used for a hardware-normalized comparison.

The available mobile platform in Figure 3 illustrates how the observation and goal interfaces could connect to physical equipment. Its ROS 1 stack, ZED depth/pose output, GNSS/IMU configuration, and goal-based autonomous navigation are supplied by the platform owner. The photograph documents hardware availability; all performance results in this article come from simulation.

\begin{center}
\begin{minipage}{\linewidth}
\centering
\includegraphics[width=\linewidth,height=.72\textheight,keepaspectratio]{docs/thesis/figures/robot_platform_20260928/robot_platform.pdf}
\captionof{figure}{Available laboratory mobile platform and proposed sensor-to-planner interfaces. Numbers identify the stereo camera, laser scanner, positioning equipment, and mobile base. The interface diagram connects observation goals to ROS 1 navigation; reported performance is measured in simulation.}\label{fig:writing-3}
\end{minipage}
\end{center}

\hypertarget{conditional-benefit-of-category-information}{%
\subsection{Conditional Benefit of Category Information}\label{conditional-benefit-of-category-information}}

\par\noindent\begin{minipage}{\linewidth}
\textbf{Table 2. Complete paired confirmation results. Relative change is the difference between means divided by the G mean.}

\begin{center}
\small
\setlength{\tabcolsep}{4pt}
\setlength{\ReviewTableWidth}{\dimexpr\linewidth-6\tabcolsep\relax}
\begin{tabular}{@{}>{\raggedright\arraybackslash}p{0.25\ReviewTableWidth} >{\raggedright\arraybackslash}p{0.25\ReviewTableWidth} >{\raggedright\arraybackslash}p{0.25\ReviewTableWidth} >{\raggedright\arraybackslash}p{0.25\ReviewTableWidth}@{}}
\toprule
Parent / configuration & G: \(J_5\) & S: \(J_5\) & S−G \\
\midrule
P00 / \(h_0\) & 0.772200 & 0.772200 & 0 \\
P00 / \(h_1\) & 0.672703 & 0.760456 & +0.087753 \\
P01 / \(h_0\) & 0.727281 & 0.727281 & 0 \\
P01 / \(h_1\) & 0.638118 & 0.723586 & +0.085468 \\
Equal-parent mean & 0.702576 & 0.745881 & +0.043305 (+6.1638\%) \\
\bottomrule
\end{tabular}
\end{center}
\end{minipage}


All runs use 42 actions, avoid collisions, and return exactly to the initial pose. Both methods attain mean coverage of 0.952544. Mean surface F1 increases from 0.737490 to 0.782951, so the joint improvement comes from surface quality, without an accompanying coverage or localization improvement. In both \(h_1\) configurations, the methods choose different turns at action 19 after the same nonsemantic history. S raises surface recall from 0.5958 to 0.7249 in P00 and from 0.5555 to 0.6845 in P01. Both \(h_0\) pairs follow identical trajectories and obtain identical scores.

The gain persists descriptively at 2 and 10 cm (Table 3). These are alternative evaluations of the same four conditions, not additional trials. With only two previously seen parent layouts, the results do not establish statistical significance or generalization to unseen geometry.

\par\noindent\begin{minipage}{\linewidth}
\textbf{Table 3. Distance-threshold sensitivity of saved confirmation results. Each row evaluates the same two parents and four conditions.}

\begin{center}
\small
\setlength{\tabcolsep}{4pt}
\setlength{\ReviewTableWidth}{\dimexpr\linewidth-6\tabcolsep\relax}
\begin{tabular}{@{}>{\raggedright\arraybackslash}p{0.25\ReviewTableWidth} >{\raggedright\arraybackslash}p{0.25\ReviewTableWidth} >{\raggedright\arraybackslash}p{0.25\ReviewTableWidth} >{\raggedright\arraybackslash}p{0.25\ReviewTableWidth}@{}}
\toprule
Distance threshold & G mean joint score & S mean joint score & Relative mean change \\
\midrule
2 cm & 0.641180 & 0.676470 & +5.5040\% \\
5 cm (primary) & 0.702576 & 0.745881 & +6.1638\% \\
10 cm & 0.723452 & 0.768385 & +6.2110\% \\
\bottomrule
\end{tabular}
\end{center}
\end{minipage}


\begin{center}
\begin{minipage}{\linewidth}
\centering
\includegraphics[width=\linewidth,height=.72\textheight,keepaspectratio]{docs/thesis/figures/virtual_method_20260928/publication_main.pdf}
\captionof{figure}{(A) All G/S pairs across two parent layouts and two configurations, including both ties. (B) Condition-wise comparison of four CPU methods. (C, D) Coverage and surface F1; points denote the four conditions and horizontal caps indicate means, without confidence intervals. G/S in B reuse A and provide no additional independent evidence. SWAP-I and VISTA-I are mechanism adaptations, not complete-system rankings.}\label{fig:writing-4}
\end{minipage}
\end{center}

Figures 5 and 6 connect the aggregate surface scores to the saved geometry. All four conditions use a paired display: G and S share camera, scale, and lighting within each condition. The two identical \(h_0\) meshes remain visible as controls. Local magnification explains which side surfaces differ without introducing a selected local quality metric.

\begin{center}
\begin{minipage}{\linewidth}
\centering
\includegraphics[width=\linewidth,height=.72\textheight,keepaspectratio]{docs/thesis/figures/scene_details_20260928/mesh_comparison.pdf}
\captionof{figure}{Saved final public-ROI TSDF meshes for all four conditions. Each G/S pair shares the same orthographic camera, scale, neutral material and lighting; all saved triangles are projected. h0 and h1 are viewed from the rear-left and rear-right respectively. F1 and J5 are copied from frozen measurements. Dashed windows identify the subsequent detail views. Both h0 pairs are array-identical. Shading is illumination, not error; no completion, smoothing or ground-truth surface is added.}\label{fig:writing-5}
\end{minipage}
\end{center}

\begin{center}
\begin{minipage}{\linewidth}
\centering
\includegraphics[width=\linewidth,height=.72\textheight,keepaspectratio]{docs/thesis/figures/scene_details_20260928/mesh_details.pdf}
\captionof{figure}{Magnified windows of the identical mesh projections. Windows are determined from the declared lateral-rib geometry and use a fixed 3.40 m by 2.55 m projected extent across conditions. These are display crops of saved surfaces, not new reconstructions. The h1 comparisons show the additional observed side surfaces; both h0 pairs retain identical geometry. No local quality score is introduced.}\label{fig:writing-6}
\end{minipage}
\end{center}

\hypertarget{correction-of-incorrect-priors}{%
\subsection{Correction of Incorrect Priors}\label{correction-of-incorrect-priors}}

The incorrect-prior development experiment uses two configurations of one parent layout. X exchanges category interpretations; Xnf further disables measured correction and anticipation of future correction. Mean \(J_5\) increases from 0.655065 for Xnf to 0.672254 for X, a relative improvement of 2.6241\%. This estimates the combined correction policy. Same-state posterior interventions support an effect of measured updates on the next action, but provide no additional counterfactual endpoint trajectories. At the stricter 2 cm threshold, X−Xnf is −0.0021593573 for \(h_1\), demonstrating that correction does not improve every configuration at every threshold.

\begin{center}
\begin{minipage}{\linewidth}
\centering
\includegraphics[width=\linewidth,height=.72\textheight,keepaspectratio]{docs/thesis/figures/virtual_method_20260928/publication_ablation.pdf}
\captionof{figure}{Endpoints for four policies in the two P00 development configurations, with X−Xnf differences at 2, 5, and 10 cm. Xnf disables measured correction and anticipated future correction together; it is not a single-posterior-update ablation. This development batch is kept separate from the held-out-noise confirmation in Figure 4.}\label{fig:writing-7}
\end{minipage}
\end{center}

\begin{center}
\begin{minipage}{\linewidth}
\centering
\includegraphics[width=\linewidth,height=.72\textheight,keepaspectratio]{docs/thesis/figures/virtual_method_20260928/publication_timeline.pdf}
\captionof{figure}{Saved P00/\(h_1\) development logs. (a) Category information establishes a prior within the shared prefix; the decision after step 18 selects budgeted action 19. (b) Geometry observed at step 28 corrects the misleading prior and changes the next action. Curves show uncalibrated configuration weights. This mechanism illustration does not increase the number of endpoint samples.}\label{fig:writing-8}
\end{minipage}
\end{center}

\hypertarget{external-mechanisms-and-computational-cost}{%
\subsection{External Mechanisms and Computational Cost}\label{external-mechanisms-and-computational-cost}}

The common CPU environment includes SWAP-I inspection and VISTA-I directional-coverage adaptations. Each completes four valid tasks, with mean \(J_5\) of 0.572223 and 0.640977, respectively, compared with 0.745881 for S. This comparison shows that observation allocation matters in the present setting. Because objectives and planning details also change, the differences cannot be attributed entirely to category information or interpreted as rankings of the original complete systems.

A native TARE node separately completes four migration tasks {[}5{]}. Narrow-field-of-view conversion, waypoint projection, budgeted waiting, and return behavior depend on the adapter, so these runs serve as engineering transfer evidence rather than a matched performance baseline.

\par\noindent\begin{minipage}{\linewidth}
\textbf{Table 4. Four-condition means for CPU mechanism adaptations. G/S reuse Table 2 trajectories. All methods have mean coverage 0.952544. Runtime covers the planner and observation record only, excluding sensing, TSDF fusion, and endpoint evaluation.}

\begin{center}
\small
\setlength{\tabcolsep}{4pt}
\setlength{\ReviewTableWidth}{\dimexpr\linewidth-6\tabcolsep\relax}
\begin{tabular}{@{}>{\raggedright\arraybackslash}p{0.25\ReviewTableWidth} >{\raggedright\arraybackslash}p{0.25\ReviewTableWidth} >{\raggedright\arraybackslash}p{0.25\ReviewTableWidth} >{\raggedright\arraybackslash}p{0.25\ReviewTableWidth}@{}}
\toprule
Method & Surface F1 at 5 cm & \(J_5\) & Planner and record time per task (s) \\
\midrule
G & 0.737490 & 0.702576 & 0.067853 \\
S & 0.782951 & 0.745881 & 0.123555 \\
SWAP-I & 0.600887 & 0.572223 & 1.256934 \\
VISTA-I & 0.672911 & 0.640977 & 1.388679 \\
\bottomrule
\end{tabular}
\end{center}
\end{minipage}


\hypertarget{extended-budget-noise-and-geometry-tests}{%
\subsection{Extended Budget, Noise, and Geometry Tests}\label{extended-budget-noise-and-geometry-tests}}

A fixed extension contains 128 trials with the same controller, prior settings, 18-action prefix, and evaluation. The original P00/P01 layouts, two configurations, budgets 30/42/54, two seeds, and four policies define 96 trials; two same-family geometry variants at budget 42 define another 32. Seeds are 2026092801 and 2026092802. One P00 rib is shifted by 0.25 m along local y; three P01 ribs have their upper boundary lowered from 1.80 to 1.35 m. Each transformation is applied to both public hypotheses, sensed geometry, and evaluation references, preserving the safe graph and prefix. Variants are not unseen natural equipment, and noise repeats are not independent layouts.

All 128 declared trials were attempted once without replacements or retries. There are 96 qualified completions with no collisions and exact returns, and 32 planning-feasibility failures at budget 30. After the shared prefix, the controller could not find a return plan satisfying its 80\% predicted-coverage gate for both templates. This is a boundary of the current planning constraints, not proof that any algorithm must fail in the physical task. All 19 observed frames and failure records are retained. Offline scores for partial failure maps are reported separately and are not successful endpoints.

\par\noindent\begin{minipage}{\linewidth}
\textbf{Table 5. Complete J5 means in the new batch. Budget 30 has 0/32 qualified trials; each other row has 32/32, with eight trials per policy. Failed groups remain blank.}

\begin{center}
\small
\setlength{\tabcolsep}{4pt}
\setlength{\ReviewTableWidth}{\dimexpr\linewidth-10\tabcolsep\relax}
\begin{tabular}{@{}>{\raggedright\arraybackslash}p{0.17\ReviewTableWidth} >{\raggedright\arraybackslash}p{0.17\ReviewTableWidth} >{\raggedright\arraybackslash}p{0.17\ReviewTableWidth} >{\raggedright\arraybackslash}p{0.17\ReviewTableWidth} >{\raggedright\arraybackslash}p{0.16\ReviewTableWidth} >{\raggedright\arraybackslash}p{0.16\ReviewTableWidth}@{}}
\toprule
Condition / budget & G & S & X & Xnf & S/G change \\
\midrule
Original / 30 & --- & --- & --- & --- & Not compared \\
Original / 42 & 0.702648 & 0.746567 & 0.652539 & 0.641425 & +6.2505\% \\
Original / 54 & 0.926298 & 0.919819 & 0.931050 & 0.908197 & −0.6995\% \\
Variants / 42 & 0.704562 & 0.734357 & 0.664307 & 0.653192 & +4.2289\% \\
\bottomrule
\end{tabular}
\end{center}
\end{minipage}


At budget 42, S/G has four wins and four ties, giving a 6.2505\% mean improvement at identical mean coverage of 0.952544. At budget 54, it has four ties and four losses, with a −0.6995\% relative mean difference. Coverage reaches 1.0 for both policies; G/S surface recall is 0.985764/0.983826 and precision is 0.873764/0.863789. The negative difference mainly accompanies lower precision, consistent with a remaining gap between template exposure and measured TSDF quality. This decomposition does not isolate a causal fusion-error mechanism. X also exceeds G/S at this budget, so the finite-lookahead exposure objective does not guarantee that a correct prior gives the best measured surface in every regime.

\begin{center}
\begin{minipage}{\linewidth}
\centering
\includegraphics[width=\linewidth,height=.72\textheight,keepaspectratio]{docs/thesis/figures/expansion_20260928/budget_effects.pdf}
\captionof{figure}{Budget sensitivity on the original layouts: completed endpoint quality, all paired contrasts, and qualification counts. No completed-performance mean is reported at B=30 because no endpoint qualifies. Points represent matched layout, configuration, and noise conditions; short bars denote complete-pair means, not confidence intervals over independent scenes.}\label{fig:writing-9}
\end{minipage}
\end{center}

The rib-shift and lower-rib variants yield mean S/G improvements of 5.0337\% and 3.3974\%, respectively, or 4.2289\% together, with four wins and four ties. At 2/5/10 cm, the original-layout improvements at budget 42 are 5.4991\%/6.2505\%/6.3063\%, and variant improvements are 3.6126\%/4.2289\%/4.3803\%. Original-layout differences at budget 54 are −1.2217\%/−0.6995\%/−0.2866\%. Thus, the budget boundary is not a consequence of reporting only the 5 cm threshold.

The whole correction policy X/Xnf improves mean J5 by 1.7328\% and 2.5163\% on original layouts at budgets 42 and 54, and by 1.7016\% on the variants at budget 42. At budget 42, P00 and its variant improve while P01 and its variant tie, showing that correction benefits also depend on the condition. Full per-trial precision, recall, F1, coverage, and qualification remain available with the paired comparisons.

\begin{center}
\begin{minipage}{\linewidth}
\centering
\includegraphics[width=\linewidth,height=.72\textheight,keepaspectratio]{docs/thesis/figures/expansion_20260928/scene_and_threshold_effects.pdf}
\captionof{figure}{Geometry perturbations and threshold sensitivity at B=42. The upper panel retains every paired result for the two original layouts and two same-family variants; the lower panel reports signed means at 2, 5, and 10 cm. Open and filled markers denote h0 and h1. Variants do not establish unseen-category generalization.}\label{fig:writing-10}
\end{minipage}
\end{center}

\hypertarget{executed-costs-and-reconstruction-progress}{%
\subsection{Executed Costs and Reconstruction Progress}\label{executed-costs-and-reconstruction-progress}}

The saved primitive-action strips in Figure 1 distinguish distance from sensing and turning cost. P00/h1 uses 28 m and 14 turns for G versus 26 m and 16 turns for S, with 42 paid actions in each case. P01/h1 uses 28 m and 14 turns for both policies, but their action sequences differ. Every paid primitive obtains RGB-D and a planar scan; these local experiments use no additional observe primitive. The measured category gain therefore does not require a longer route or more sensor frames.

To examine how reconstruction gains develop, we reintegrated all eight original trajectories at 18, 24, 30, 36, and 42 paid actions, retaining their recorded observations, poses, TSDF settings, and surface evaluator. All 16 previously available prefix/endpoint checkpoints reproduce the original geometry and metric values, with zero metric difference. Planar coverage is identical within every pair at all five checkpoints and reaches its endpoint value by action 30. In both h1 conditions, later quality differences are associated with increased surface recall and F1. P00/h1 nevertheless exhibits a joint-score difference of −0.000819 at action 24 before becoming positive at the later checkpoints; both h0 pairs remain tied throughout. These process measurements reuse the original trajectories and add no independent trials. Intermediate scores do not imply successful terminal return.

\begin{center}
\begin{minipage}{\linewidth}
\centering
\includegraphics[width=\linewidth,height=.72\textheight,keepaspectratio]{docs/thesis/figures/article_reconstruction_20260928/checkpoint_quality_components.pdf}
\captionof{figure}{Offline reintegration and measurement of all eight saved V36 trajectories at 18, 24, 30, 36 and 42 paid actions. Columns separate actual planar coverage, 5 cm precision, recall, F1 and joint J5; rows retain all four paired conditions. The action axis is accompanied by actual cumulative translation distances. Markers are measured checkpoints and connecting segments only visual guides. Intermediate stages are not successful terminal episodes. Prefix and final stages are checked against original saved outputs; 40 checkpoints do not constitute additional independent trajectories.}\label{fig:writing-11}
\end{minipage}
\end{center}

\begin{center}
\begin{minipage}{\linewidth}
\centering
\includegraphics[width=\linewidth,height=.72\textheight,keepaspectratio]{docs/thesis/figures/article_reconstruction_20260928/checkpoint_actual_meshes.pdf}
\captionof{figure}{Actual reintegrated TSDF meshes at the fixed display checkpoints 18, 30 and 42 for all four G/S pairs. All extracted triangles are rendered; each configuration shares one camera, metric scale, material and lighting across methods and time. The h0/h1 views use predefined rear-left/rear-right cameras. No hole completion, smoothing or template surface is added. F1 and J5 come from the offline measurements; shading is illumination, not error. Both coincident h0 conditions remain visible.}\label{fig:writing-12}
\end{minipage}
\end{center}

\hypertarget{scene-level-development-and-information-timing}{%
\subsection{Scene-Level Development and Information Timing}\label{scene-level-development-and-information-timing}}

The complete development design comprises AISLE, CELL, and LOOP layouts, four methods, one common seed, and a 160-action budget: 12 attempted episodes. Eleven complete the original execution-and-evaluation pipeline. CELL/G completes collision-free motion and full-pose return but its evaluator rejects one finite triangle of area \(4.24\times10^{-13}\,\mathrm{m}^2\). This remains an original evaluation failure, not a navigation failure or a zero-quality result.

\par\noindent\begin{minipage}{\linewidth}
\textbf{Table 6. Original scene-level development endpoints. Entries are \(Q/J_{\rm nav}\); a dash denotes unavailable original evaluation. Each entry is one episode.}

\begin{center}
\small
\setlength{\tabcolsep}{4pt}
\setlength{\ReviewTableWidth}{\dimexpr\linewidth-8\tabcolsep\relax}
\begin{tabular}{@{}>{\raggedright\arraybackslash}p{0.2\ReviewTableWidth} >{\raggedright\arraybackslash}p{0.2\ReviewTableWidth} >{\raggedright\arraybackslash}p{0.2\ReviewTableWidth} >{\raggedright\arraybackslash}p{0.2\ReviewTableWidth} >{\raggedright\arraybackslash}p{0.2\ReviewTableWidth}@{}}
\toprule
Layout & NBV & G & B & S \\
\midrule
AISLE & 0.857/0.857 & 0.857/0.857 & 0.857/0.857 & 0.857/0.857 \\
CELL & 0.736/0.690 & ---/--- & 0.819/0.767 & 0.819/0.767 \\
LOOP & 0.649/0.621 & 0.649/0.621 & 0.649/0.621 & 0.649/0.621 \\
\bottomrule
\end{tabular}
\end{center}
\end{minipage}


A common numerical supplement applies the same face-area threshold, \(5\times10^{-13}\,\mathrm{m}^2\), to all 12 predictions. Eleven unchanged-input evaluations are verified and reused; only CELL/G requires a new derived surface measurement, with \(C_{\rm nav}=0.936450\), \(Q=0.817264\), and \(J_{\rm nav}=0.765328\). No new route or TSDF is generated. B/S exceed this G score by only 0.001193, approximately 0.156\%; S and B tie in all three layouts. The supplement preserves original failure status and does not establish a sharing benefit.

\begin{center}
\begin{minipage}{\linewidth}
\centering
\includegraphics[width=\linewidth,height=.72\textheight,keepaspectratio]{docs/thesis/figures/article_online_20260928/development12_turnmarkers/online_routes_01.pdf}
\captionof{figure}{All 12 declared slots at a common scale. Circles mark turns, squares paid observations, and arrows saved headings. The CELL/G blank denotes original evaluation failure despite a completed saved route. Background geometry is for offline explanation.}\label{fig:writing-13}
\end{minipage}
\end{center}

\begin{center}
\begin{minipage}{\linewidth}
\centering
\includegraphics[width=\linewidth,height=.72\textheight,keepaspectratio]{docs/thesis/figures/article_online_20260928/development12_turnmarkers/online_quality_and_cost.pdf}
\captionof{figure}{Original qualified endpoints and recorded costs; crosses denote unavailable entries, not zeros. Planning times exclude offline evaluation and reflect concurrent workload. This development comparison does not estimate performance on unseen layouts.}\label{fig:writing-14}
\end{minipage}
\end{center}

AISLE's four methods have identical executed actions, poses, and mesh arrays: 24 m, 54 turns, 10 explicit observations, and 161 frames each. The first informative same-category peer appears at action 136, after return begins at action 126. Inspection identified depth support connected through the floor, delaying reliable instance association. A common geometric revision excludes a measured floor band only from association candidates, using current depth and calibrated mounting height. Support, height, and noise checks control acceptance; rejected frames retain the original input. All methods receive the same rule, and full depth remains available to occupancy and TSDF. This assumes a static horizontal floor and known camera geometry.

Replaying the same 161 saved frames reproduces the original evidence exactly. With the revision, first applied feedback moves from action 85 to 8, and first informative peer/S--B posterior difference from 136 to 82. Uncertain instance-frame records decrease from 278 to 3. However, applied feedback decreases from 21 to 18, and the second cabinet's first own feedback shifts from 85 to 87. These results establish earlier information availability in a fixed-path replay, not improved association accuracy against ground truth, a changed route, or higher reconstruction quality. The common-frontend closed-loop comparison is reported separately in Section 5.8; the fixed-path replay itself establishes no new planning benefit.

\begin{center}
\begin{minipage}{\linewidth}
\centering
\includegraphics[width=\linewidth,height=.72\textheight,keepaspectratio]{docs/thesis/figures/article_scene_meshes_20260928/aisle_gbs_v2/aisle_gt_g_s_preview.pdf}
\captionof{figure}{Ground truth and saved G/S meshes share view and scale; all 205,566 predicted triangles are rendered with overlaid XY routes. Colors encode height and lighting, not error. G/S geometry is identical. GT includes unobservable surfaces beyond the fixed scoring reference.}\label{fig:writing-15}
\end{minipage}
\end{center}

Near-complete AISLE navigable coverage (0.999289) coexists with macro F1 0.857363 and second-facility recall 0.557492. Figure 15 connects route coverage to incomplete equipment surfaces without depicting an absent G/S advantage. Together, these results distinguish the local category benefit from the additional scene-level requirements of stable association and timely, actionable information.

\hypertarget{closed-loop-association-ablation}{%
\subsection{Closed-Loop Association Ablation}\label{closed-loop-association-ablation}}

A further 12 episodes apply the measured-ground association rule to every method on the same three development layouts, budget, and noise realization. All complete collision-free motion, full-pose return, and the common numerical evaluation. Table 7 compares the original predictions under that same evaluator with the new closed-loop predictions. This is a common frontend ablation; no historical route is regenerated with a different controller.

\par\noindent\begin{minipage}{\linewidth}
\textbf{Table 7. Joint score before and after the common association correction. Each cell is original → ground-aware; † denotes the separately derived score after an original evaluation failure.}

\begin{center}
\small
\setlength{\tabcolsep}{4pt}
\setlength{\ReviewTableWidth}{\dimexpr\linewidth-8\tabcolsep\relax}
\begin{tabular}{@{}>{\raggedright\arraybackslash}p{0.2\ReviewTableWidth} >{\raggedright\arraybackslash}p{0.2\ReviewTableWidth} >{\raggedright\arraybackslash}p{0.2\ReviewTableWidth} >{\raggedright\arraybackslash}p{0.2\ReviewTableWidth} >{\raggedright\arraybackslash}p{0.2\ReviewTableWidth}@{}}
\toprule
Layout & NBV & G & B & S \\
\midrule
AISLE & 0.857 →\allowbreak{} 0.653 & 0.857 →\allowbreak{} 0.653 & 0.857 →\allowbreak{} 0.653 & 0.857 →\allowbreak{} 0.653 \\
CELL & 0.690 →\allowbreak{} 0.465 & 0.765† →\allowbreak{} 0.534 & 0.767 →\allowbreak{} 0.531 & 0.767 →\allowbreak{} 0.531 \\
LOOP & 0.621 →\allowbreak{} 0.730 & 0.621 →\allowbreak{} 0.730 & 0.621 →\allowbreak{} 0.730 & 0.621 →\allowbreak{} 0.730 \\
\bottomrule
\end{tabular}
\end{center}
\end{minipage}


The correction improves LOOP's joint score by 0.109889 but lowers it in AISLE and CELL. Earlier usable information therefore does not guarantee a better observation policy. AISLE's route first changes at action 6: translation decreases from 24 to 22 m, while turns increase from 54 to 56 and explicit observations from 10 to 16. Coverage and macro F1 both decrease. In LOOP, translation remains 20 m, but changed observation allocation improves both quantities. The full paired coverage/F1/product figure and per-instance records accompany the data.

\begin{center}
\begin{minipage}{\linewidth}
\centering
\includegraphics[width=\linewidth,height=.72\textheight,keepaspectratio]{docs/thesis/figures/article_ground_20260928/fixed_g_print_v1/ground_fixed_g_routes.pdf}
\captionof{figure}{Fixed geometric method G in all three development layouts, comparing the common association component at one metric scale. Dashed and solid lines are actual original and ground-aware trajectories; circles mark turns, squares extra paid observations, and arrows recorded headings. CELL retains its original evaluation failure and separately derived score. The full four-method grid is available with the data. Equipment outlines provide offline context.}\label{fig:writing-16}
\end{minipage}
\end{center}

Within each version, S and B execute identical actions and obtain identical quality in all three layouts. Under the correction, their posteriors differ in 23 instance--step records in CELL and 93 in LOOP, but no compared candidate score or executed action changes. A separate analysis of 138 category-conditioned direct-view pools from the original study and the first AISLE correction pair also excludes a ranking change over the one-peer reliability range, even with relaxed peer availability. These dependent decision records are mechanism diagnostics, not additional independent trials; they do not cover diagnostic alternatives or future candidate pools.

Saved AISLE decisions reveal another limitation: nominal scales displaced from measured supports can retain positive unseen-area rewards inside previously acquired view cones. A subsequent, separately frozen version tests common same-center exposure deduplication. It excludes previously attempted nominal exposure from predicted new-area eligibility while leaving measured supports and TSDF unchanged. Repeated-view precision benefits and missing-depth recovery remain unmodeled. The original scene-level, ground-aware, and exposure-deduplicated versions have separate configurations and results; Table 7 contains only the first two.

\begin{center}
\begin{minipage}{\linewidth}
\centering
\includegraphics[width=\linewidth,height=.72\textheight,keepaspectratio]{docs/thesis/figures/article_ground_meshes_20260928_v2/ground_frontend_meshes_G.pdf}
\captionof{figure}{Fixed geometric method G, all three development layouts, and the two association versions. Reference geometry and complete saved meshes share an orthographic camera, world scale, and height coloring. All 1,055,521 stored triangles across nine panels are retained. Color is not reconstruction error. CELL's original failed evaluation and its separately derived score remain distinguished. Full-frame fusion can reconstruct surfaces of instances not registered by the semantic frontend.}\label{fig:writing-17}
\end{minipage}
\end{center}

\hypertarget{exposure-discounting-and-executed-mechanisms}{%
\subsection{Exposure Discounting and Executed Mechanisms}\label{exposure-discounting-and-executed-mechanisms}}

The final 12-episode ablation retains the measured-ground frontend and discounts previously exposed nominal opportunities at the same optical center. All three development layouts and all four methods complete collision-free full-pose return under the same budget and seed. Table 8 retains every method and layout; the original, ground-aware, and exposure-discounted batches together contain 36 attempts, not 36 independent scenes. The original CELL/G evaluation failure remains recorded separately from its derived common numerical score.

\par\noindent\begin{minipage}{\linewidth}
\textbf{Table 8. Joint score with ground-aware association, before → after exposure discounting. Each value is one reviewed endpoint; all 12 new episodes qualify.}

\begin{center}
\small
\setlength{\tabcolsep}{4pt}
\setlength{\ReviewTableWidth}{\dimexpr\linewidth-8\tabcolsep\relax}
\begin{tabular}{@{}>{\raggedright\arraybackslash}p{0.2\ReviewTableWidth} >{\raggedright\arraybackslash}p{0.2\ReviewTableWidth} >{\raggedright\arraybackslash}p{0.2\ReviewTableWidth} >{\raggedright\arraybackslash}p{0.2\ReviewTableWidth} >{\raggedright\arraybackslash}p{0.2\ReviewTableWidth}@{}}
\toprule
Layout & NBV & G & B & S \\
\midrule
AISLE & 0.653 →\allowbreak{} 0.783 & 0.653 →\allowbreak{} 0.783 & 0.653 →\allowbreak{} 0.734 & 0.653 →\allowbreak{} 0.734 \\
CELL & 0.465 →\allowbreak{} 0.515 & 0.534 →\allowbreak{} 0.787 & 0.531 →\allowbreak{} 0.787 & 0.531 →\allowbreak{} 0.787 \\
LOOP & 0.730 →\allowbreak{} 0.736 & 0.730 →\allowbreak{} 0.736 & 0.730 →\allowbreak{} 0.690 & 0.730 →\allowbreak{} 0.690 \\
\bottomrule
\end{tabular}
\end{center}
\end{minipage}


Exposure discounting improves AISLE and CELL scores for every method. In LOOP, G/NBV improve slightly while B/S decrease by 0.040235. It therefore addresses a repeated-opportunity problem without providing a general reconstruction guarantee. Moreover, all AISLE scores remain below the original version's 0.856753. The component revisions cannot be presented as a sequence of uniformly improving systems.

Two saved causal sequences distinguish semantic influence from predicted information value. In AISLE, B and G receive identical sensor arrays through frame 31 and apply the same geometric evidence. At decision 27, the category prior changes the selected direct-view target. The first actual divergence is an opposite turn at paid action 32; the selected target observations complete at actions 41 for B and 44 for G. Both pass budget and execution checks, yet B's final joint score is lower by 0.048716. The cue demonstrably affects executed decisions in this case without improving quality.

In CELL, G and NBV have identical observations through frame 19 and identical scores for all 32 direct candidates at decision 17. G additionally considers seven diagnostic options. Its selected diagnostic score is 0.082740, including an anticipated information increment of 0.006106, above the best direct score of 0.078458. At action 20, G observes while NBV turns; their selected macros finish at actions 20 and 29. However, G's diagnostic frame applies no new geometric feedback, leaves its uniform structure posterior unchanged, and adds no known planar cells. Its final score exceeds NBV by 0.271044 after the complete subsequent trajectories. This is a diagnostic-planning component case, not evidence that the anticipated information was obtained or that this single frame caused the entire endpoint improvement.

\begin{center}
\begin{minipage}{\linewidth}
\centering
\includegraphics[width=\linewidth,height=.72\textheight,keepaspectratio]{docs/thesis/figures/article_exposure_20260928/complete12/exposure_paired_quality.pdf}
\captionof{figure}{Coverage, full four-instance macro F1, and their product for every ground-aware and exposure-discounted pair. Markers denote actual terminal measurements from three development layouts; neither methods nor frames add independent layout replicates. Both revisions retain full-depth fusion and the same evaluation. The complete route grid accompanies the data.}\label{fig:writing-18}
\end{minipage}
\end{center}

\begin{center}
\begin{minipage}{\linewidth}
\centering
\includegraphics[width=\linewidth,height=.72\textheight,keepaspectratio]{docs/thesis/figures/article_mechanism_cases_20260928_final/mechanism_cases.pdf}
\captionof{figure}{Two development examples at the same metric scale, retaining complete executed paths and return. D denotes target selection, F the first actual action difference, and O completion of a paid target view. AISLE compares category-conditioned B with geometric G; CELL compares geometric lookahead G with direct NBV. Arrows show saved camera headings. The signed endpoint differences include all later actions; they are not per-view effects. CELL's diagnostic frame does not provide fresh feedback. Neither panel tests the S/B sharing increment.}\label{fig:writing-19}
\end{minipage}
\end{center}

The three exposure-discounted S/B pairs also have identical complete 161-frame actual sensor sequences and actions. In CELL, the first posterior difference occurs at step 73 and the first candidate-score difference at step 100, without changing execution. In LOOP, posteriors differ from step 65 but candidate scores remain unchanged; AISLE has neither difference. Thus none of the nine S/B pairs across the three versions supplies an executed sharing effect. The predeclared requirement for at least one such causal action difference is not met, so the reserved 48-run unseen-layout matrix is not executed. Unstarted layouts are not failures or zero-quality measurements.

\hypertarget{discussion-and-applicability}{%
\section{Discussion and Applicability}\label{discussion-and-applicability}}

Category information is useful here because it can select an observation side while geometry remains ambiguous. When both methods choose the same route, the observed gain is zero. In particular, G resolves an equal-valued left/right decision by choosing left, which coincides with S in both original \(h_0\) configurations at budget 42. The extension preserves this budget's gain under new noise and same-family geometry, while a larger budget changes policy ordering. Ambiguity strength and natural category reliability have not been independently manipulated or calibrated. Public templates, synthetic category markers, exact poses, a shared safe graph, and a fixed prefix define these local experiments. Their region- and threshold-dependent surface evaluation limits the conclusions to local facility documentation. The separate full-mesh extension in Section 5.7 retains the navigation and perception assumptions and does not establish a shared-semantic performance advantage.

The analytical model identifies conditions for a reduced or absent semantic advantage; it does not predict the observed high-budget sign reversal. Its optimal decisions use a correctly specified utility and observation model, whereas the executed planner uses finite lookahead and uncalibrated exposure proxies. The measured reversal therefore remains an empirical boundary requiring reconstruction-level evidence, not a consequence established by the analytical example.

An earlier six-parent shared-reliability development study, distinct from the new three-layout extension in Section 5.7, provides an additional boundary. Under a different controller and \(J_{\rm nav}\) task, the nominal category--structure condition gives five S/G ties and one loss across six parent layouts, with a relative mean difference of −0.3301\%. Under a mismatched relation, all six comparisons between S and the ordinary Bayesian semantic baseline B tie. Thus, implemented observation, belief update, and replanning do not necessarily produce an endpoint improvement. One nominal B control is missing because of interruption, and the full development matrix remains incomplete. These results cannot be pooled with the present \(J_5\) results and do not establish an effective shared-reliability contribution. They reinforce the narrower interpretation: the demonstrated benefit concerns category information in the original controlled setting, rather than a general advantage of increasingly complex semantic mechanisms.

\hypertarget{conclusion}{%
\section{Conclusion}\label{conclusion}}

Motivated by active mapping needs in industrial facilities, this study implements category-prior and geometric-feedback active observation under a motion budget. The original confirmation yields a 6.1638\% mean category benefit, and the separate incorrect-prior study supports a 2.6241\% whole-policy correction benefit. A fixed 128-trial extension reproduces a 6.2505\% category improvement at budget 42 and a 4.2289\% improvement on same-family variants, while retaining low-budget infeasibility and a high-budget reversal. The evidence supports a conditional role for public facility knowledge in allocating limited observations. Across three scene-level versions, S/B quality ties persist on all three development layouts. The exposure-discounted version further demonstrates that a category cue can change an executed target yet reduce final quality. The executable multi-instance extension has therefore not established a shared-reliability benefit. Future work should improve agreement between predicted observation value and measured surface quality before extending map maintenance across tasks and physical deployment.

\hypertarget{reproducibility-and-data}{%
\section*{Reproducibility and Data}\label{reproducibility-and-data}}

The NSO research repository preserves scientific outputs, original hashes, and restored historical source archives. The separate cleaned SGAM snapshot is a release view rather than a substitute for every historical dependency. Exact reproduction should use the source and input manifests of the relevant batch; a successful source audit does not imply that every legacy runtime dependency is currently installed. New acquisitions seal their own implementation and retain separate records.

The extension preserves the fixed configuration, executed source hashes, and all 128 declared trial records. Its \path{analysis_review/complete_endpoint_metrics.csv} contains endpoint components and qualification; \path{complete_paired_metrics.csv} matches layout, configuration, budget, and seed. The batch \texttt{source\_freeze.json} binds implementation and inputs. Plotting scripts read sealed results, with separate provenance for saved meshes and the laboratory photographs. Historical confirmation, incorrect-category, and external-mechanism batches remain separate. Software versions are given in Section 5.1; the accompanying platform reproduction note and thesis Appendix A provide dependency and verification instructions. Any fresh acquisition must use a separate output rather than overwrite archived evidence.

\hypertarget{references}{%
\section*{References}\label{references}}

\begin{enumerate}
\def\labelenumi{\arabic{enumi}.}
\item
  Chaplot D. S., Gandhi D., Gupta S., Gupta A., Salakhutdinov R. Learning to Explore using Active Neural SLAM. ICLR, 2020. \href{https://arxiv.org/abs/2004.05155}{Author paper}.
\item
  Dharmadhikari M., Alexis K. Semantics-aware Exploration and Inspection Path Planning. ICRA, 2023. \href{https://arxiv.org/abs/2303.07236}{Author paper}.
\item
  Nagami K., Chen T., Yu J., Shorinwa O., Adang M., Dougherty C., Cristofalo E., Schwager M. VISTA: Open-Vocabulary, Task-Relevant Robot Exploration With Online Semantic Gaussian Splatting. IEEE Robotics and Automation Letters, 11(3):3150--3157, 2026. \href{https://doi.org/10.1109/LRA.2026.3653276}{doi:10.1109/LRA.2026.3653276}.
\item
  Zhou Q.-Y., Park J., Koltun V. Open3D: A Modern Library for 3D Data Processing. arXiv:1801.09847, 2018. \href{https://doi.org/10.48550/arXiv.1801.09847}{doi:10.48550/arXiv.1801.09847}.
\item
  Cao C., Zhu H., Choset H., Zhang J. TARE: A Hierarchical Framework for Efficiently Exploring Complex 3D Environments. RSS, 2021. \href{https://doi.org/10.15607/RSS.2021.XVII.018}{doi:10.15607/RSS.2021.XVII.018}.
\item
  Chen X., Li Q., Wang T., Xue T., Pang J. GenNBV: Generalizable Next-Best-View Policy for Active 3D Reconstruction. CVPR, 2024:16436--16445. \href{https://openaccess.thecvf.com/content/CVPR2024/html/Chen_GenNBV_Generalizable_Next-Best-View_Policy_for_Active_3D_Reconstruction_CVPR_2024_paper.html}{CVF proceedings}.
\item
  Dharmadhikari M., Alexis K. Semantics-Aware Predictive Inspection Path Planning. IEEE Transactions on Field Robotics, 2025. \href{https://doi.org/10.1109/TFR.2025.3578402}{doi:10.1109/TFR.2025.3578402}.
\item
  Chen L., Zhan H., Yin H., Xu Y., Mordohai P. Understanding while Exploring: Semantics-driven Active Mapping. Advances in Neural Information Processing Systems, 38, 2025. \href{https://doi.org/10.52202/085713-1113}{doi:10.52202/085713-1113}.
\item
  World Economic Forum. Physical AI: Powering the New Age of Industrial Operations. White paper, 2025-09-04. \href{https://www.weforum.org/publications/physical-ai-powering-the-new-age-of-industrial-operations/}{Official white paper}.
\item
  Placed J. A., Strader J., Carrillo H., Atanasov N., Indelman V., Carlone L., Castellanos J. A. A Survey on Active Simultaneous Localization and Mapping: State of the Art and New Frontiers. IEEE Transactions on Robotics, 2023. \href{https://arxiv.org/abs/2207.00254}{Author paper}.
\item
  Curless B., Levoy M. A Volumetric Method for Building Complex Models from Range Images. Proceedings of SIGGRAPH, 1996:303--312. \href{https://graphics.stanford.edu/papers/volrange/}{Author project and paper}.
\item
  Kaelbling L. P., Littman M. L., Cassandra A. R. Planning and Acting in Partially Observable Stochastic Domains. Artificial Intelligence, 101(1--2):99--134, 1998. \href{https://people.csail.mit.edu/lpk/papers/aij98-pomdp.pdf}{Author paper}.
\item
  Hollinger G. A., Sukhatme G. S. Sampling-based Robotic Information Gathering Algorithms. The International Journal of Robotics Research, 33(9):1271--1287, 2014. \href{https://doi.org/10.1177/0278364914533443}{doi:10.1177/0278364914533443}.
\item
  Knapitsch A., Park J., Zhou Q.-Y., Koltun V. Tanks and Temples: Benchmarking Large-Scale Scene Reconstruction. ACM Transactions on Graphics, 36(4), 2017. \href{https://www.tanksandtemples.org/}{Author project and paper}.
\end{enumerate}
\end{document}
